\section{Bit complexity}
\label{sec:complexity}

We now bound the complete cost of the game algorithm. Its prescribed
root input is \([0,B_*]\), with unit support weights and budgets
\(k=m=n\). The comparison theorem allowed arbitrary positive rational
parameters; the bit estimates here concern the parameters reachable
from this root. All estimates hold on every random tape, including
tapes on which comparisons fail.

\begin{theorem}\label{thm:bit-complexity}
The algorithm has a uniform implementation in a standard bit-operation
model such that, for every input of bit length \(L\) and every random
tape, the running time is at most
\[
              2^{C(\log_2(L+2))^2}
\]
for an absolute constant \(C>0\). The bound includes preprocessing,
every recursive execution, intermediate arithmetic, and random-bit
generation.
\end{theorem}

\begin{proof}
There are two parts to the cost bound. We first count the solver
invocations by contracting each chain of main continuations. We then
show that the work attached to one invocation has polynomial bit
cost, including the arithmetic in boxes that extend beyond their
parents.

Write \(x=L+2\). Explicit encoding of the graph, ownership, and edge
list gives \(n=O(x)\) and \(|E|=O(x)\), and every input weight uses
at most \(L\) bits. The transformation \(t(e)=(n+1)w(e)+1\) and
the definitions of \(W,B_*\) consequently imply
\begin{equation}\label{cost:parameter-sizes}
 \log_2(B_*+1)=O(x),\quad
 g=O(x),\quad s=O(x),\quad J=O(x^{12}).
\end{equation}
For the bound on \(s\), its choice as the least sufficiently large
power of two gives \(s<2^{18}(g+1)\). These estimates use the binary
length of the weights, even when their magnitudes are exponential
in that length.

\paragraph{Counting the calls made at one rank.}
The initial parameter is \(R=n^3\). No descendant increases it, so
the largest rank on an ancestry is
\begin{equation}\label{cost:initial-rank}
 D_*=d(n^3)
     =1+\floor{\log_{10/9}(n^3)}
     =O(\log(n+1))=O(\log x).
\end{equation}
In particular, the strict rank convention gives \(D_*=1\) when
\(n=1\).

Consider one node on a main \(\Solve\) chain. Expose the calls to
\(\Solve\) inside its amplification and orientation routines, but
leave smaller-rank solver calls unexpanded. The resulting counts are
as follows.

\begin{itemize}
\item The tightening loop makes at most \(16ns+2\) calls to
\(\Ask\), each consisting of \(J\) solver invocations. Its
contribution is at most \((16ns+2)J\).
\item The \(\OrientPlus\) call performs \(J\) orientation trials.
Each trial has at most \(g+1\) levels, with at most \(2s\)
\(\Ask\) calls per level and \(J\) solver invocations per such
call. This contributes at most \(2s(g+1)J^2\).
\item The concatenated exceptional list has at most \(J(2g+2)\)
entries. Every entry triggers an amplified reweighted call of
\(J\) solver invocations, contributing at most \((2g+2)J^2\).
\end{itemize}

Proposition~\ref{prop:termination} makes all three counts unconditional.
Their sum is
\begin{equation}\label{cost:branching}
 A_0=(16ns+2)J+\bigl[2s(g+1)+2g+2\bigr]J^2.
\end{equation}
Both quadratic amplification terms are needed: orientation trials
contain amplified comparisons, and every entry returned by the
amplified orientation has its own amplified reweighted call. Entries
from minority trials and duplicate sets are all counted. Besides
these smaller-rank calls, a nonterminal solver invocation makes
exactly one unamplified main continuation, as shown in
Figure~\ref{fig:solver-recursion}.

\begin{figure}[hbp]
\centering
\begin{tikzpicture}[
  >=Latex,
  every node/.style={font=\small,align=center},
  main/.style={draw=linkblue,rounded corners=2pt,thick,
    fill=linkblue!5,minimum width=2.9cm,minimum height=1.05cm},
  reduced/.style={draw=black!65,rounded corners=2pt,
    fill=black!3,text width=3.2cm,minimum height=1.25cm},
  edge/.style={->,thick,draw=linkblue},
  side/.style={->,draw=black!65}
]
\node[main] (s0) at (0,0) {$\Solve$\\$R,\ B_0$};
\node[main] (s1) at (5.3,0) {$\Solve$\\$R,\ B_1$};
\node (dots) at (9.2,0) {$\cdots$};
\draw[edge] (s0) -- node[above,text width=2.3cm]
  {one main call\\unamplified} (s1);
\draw[edge] (s1) -- node[above,text width=2.3cm]
  {one main call\\unamplified} (dots);
\node[reduced] (r0) at (0,-3.0)
  {at most $A_0$ solver calls\\$R'\le\frac9{10}R$};
\node[reduced] (r1) at (5.3,-3.0)
  {at most $A_0$ solver calls\\$R'\le\frac9{10}R$};
\draw[side] (s0) -- node[right,text width=2.0cm]
  {amplified\\calls} (r0);
\draw[side] (s1) -- node[right,text width=2.0cm]
  {amplified\\calls} (r1);
\node[align=center,text width=12.0cm] at (4.3,-4.35)
  {Unchanged support parameters along the main chain, with
   $B_{j+1}=B_j-\lfloor B_j/(4s)\rfloor$ and at most $K$ invocations.};
\end{tikzpicture}
\caption{Contracting the main continuation chain. The lower boxes
represent all immediate smaller-rank $\Solve$ calls exposed inside
$\Ask$, $\Orient$, and $\OrientPlus$, including every trial and
repeated exceptional-set entry. The horizontal chain has only one
successor per nonterminal node and at most $K$ nodes, so it generates
at most $KA_0$ such descendants. These bounds hold on every random
tape. In the correctness proof, the same chain carries at most one
protected comparison until a local answer establishes it or the
terminal solver is reached.}
\label{fig:solver-recursion}
\end{figure}


Let \(N_d\) bound the total number of solver invocations generated
by one invocation of rank at most \(d\), including that invocation,
for the fixed graph and global parameters. A rank-zero invocation
returns directly. For positive rank, contract the main chain of at
most \(K\) invocations from Equation~\eqref{eq:main-chain-length}.
At each node of this chain there are at most \(A_0\) immediately
smaller-rank solver descendants. It follows that
\begin{equation}\label{cost:call-recurrence}
 N_0=1,\qquad N_d\le K+KA_0N_{d-1}\quad(d\ge1).
\end{equation}
There is only one main successor at each node, and its call is not
amplified. This is the reason one chain can be contracted. The
separation of one full-precision continuation from reduced-precision
calls also appears in Parys's parity-game recursion
\cite[Algorithm~2 and Section~5]{Parys2019}; here
Equation~\eqref{cost:branching} supplies the count for our particular
routines.

Equation~\eqref{cost:parameter-sizes} gives
\[
       K=O(x^2),\qquad A_0=O(x^{26}),\qquad KA_0=O(x^{28}).
\]
With \(A=KA_0\), induction in
Equation~\eqref{cost:call-recurrence} yields
\(N_d\le(K+1)(A+1)^d\). Since
\(D_*=O(\log x)\), we obtain
\begin{equation}\label{cost:call-count}
                N_{D_*}=2^{O((\log_2 x)^2)}.
\end{equation}
Every step of this count used a deterministic bound on every tape.
We must still bound the bit cost of the work done between these
solver calls.

\paragraph{Controlling integer magnitudes along an ancestry.}
The number of calls alone does not bound coordinate sizes, because
centered boxes can extend outside their parents. Consider any path
of ancestry in the full subroutine tree. Every child box is contained
in its parent except possibly a centered box
\[
        [q-A,q+A],\qquad A\in\{D,E_0\},
\]
formed by \(\Orient\). On every tape, its center \(q\) belongs to
the parent box, its radius satisfies \(A\le B/8\), and its reference
width is \(2A\le B/4\). All other transitions preserve or decrease
the reference width.

List the radii of just these centered transitions as
\(A_1,A_2,\ldots\). The first is at most \(B_*/8\). After a
transition of radius \(A_r\), the width is \(2A_r\) and cannot
increase before the next such transition. Therefore
\(A_{r+1}\le(2A_r)/8=A_r/4\). The possible outward displacement
of either endpoint is bounded by the convergent sum
\[
        \sum_{r\ge1}A_r
           \le\frac{B_*}{8}\sum_{r\ge0}4^{-r}
           =\frac{B_*}{6}.
\]
Starting from \([0,B_*]\), this proves that every actual endpoint,
center, and returned coordinate lies in
\begin{equation}\label{cost:endpoint-bound}
                       [-B_*/6,\,7B_*/6].
\end{equation}
These are real bounds on integer values; no rounding of the displayed
endpoints is intended.

The temporary sums used for midpoints, differences, scalar shifts,
and caps also have magnitude \(O(B_*)\). The operator \(F\) never
changes, and \(|t(e)|\le W\le B_*\), so each intermediate value
\(t(e)+p_j\) has the same order of magnitude. All such integers
therefore have \(O(\log_2(B_*+1))=O(x)\) bits. This includes
computations in centered boxes outside the original box, independently
of all comparison outcomes.

\paragraph{Representing the rational parameters exactly.}
Weights and budgets change only on a parameter-reducing edge into a
solver call. There are at most \(D_*\) such edges on any ancestry.
Starting from unit weights, every reachable coordinate weight is
\[
       \alpha_i=2^{-u_i^\alpha},\qquad \beta_i=2^{-u_i^\beta},
       \qquad u_i^\alpha,u_i^\beta\in\mathbb Z_{\ge0},\quad
       u_i^\alpha+u_i^\beta\le D_*.
\]
Starting from \(k=m=n\), the budgets have the forms
\begin{equation}\label{cost:budget-forms}
 k=n\frac{4^{u_k}3^{v_k}}{5^{u_k+v_k}},\qquad
 m=n\frac{4^{u_m}3^{v_m}}{5^{u_m+v_m}},\qquad
 u_k+v_k+u_m+v_m\le D_*,
\end{equation}
with nonnegative integer exponents. Hence each numerator and
denominator of a budget or weight has \(O(\log(n+1))\) bits.
One may store reduced fractions with binary integer numerators and
denominators, or store the exponents in these explicit forms.

Support sums admit a common power-of-two denominator of exponent at
most \(D_*\). The reciprocal-product sums used by \(R\) and
\(\mu\) are even simpler: for \(S\subseteq V\),
\[
       \sum_{i\in S}\frac1{\alpha_i\beta_i}
           =\sum_{i\in S}2^{u_i^\alpha+u_i^\beta}
           \le n2^{D_*}.
\]
Their bit lengths are \(O(\log(n+1))\). Exact integer arithmetic
thus computes these sums and tests \(R<1\) or \(\mu(S)>1/2\).
The support-budget tests and the height-threshold tests are likewise
exact after multiplication by positive denominators. Alternatively,
normalizing fractions with the Euclidean algorithm after each
operation gives polynomial bit cost without retaining a special
common-denominator representation.

The global integers \(g,s\) are computable by binary lengths and
integer comparisons with powers of two. Neither the rank nor the
real logarithms appearing in this proof must be evaluated by the
algorithm. The number \(2^{-J}\) is also solely an error bound;
no rational of that value is constructed.

\paragraph{Bounding the local work, including terminal iterations.}
The bit sizes are now controlled. For the small-width branch,
\(\Basic\) is used only when \(B\le128s\), and
Lemma~\ref{lem:box-solution} bounds it by
\[
                         1+128ns
\]
operator evaluations. Each evaluation scans the explicit edge list,
adding and comparing \(O(x)\)-bit integers, then clips the result.
This iteration count is independent of endpoint offsets and numerical
weight magnitudes. The other terminal branch, \(R<1\), uses no
iteration. Thus neither branch entails work proportional to \(W\)
or \(B_*\).

For a concrete representation, store a subset of \(V\) as an
\(n\)-bit incidence vector. Compute a coordinatewise median by
sorting the \(J\) entries in each coordinate and selecting entry
\((J+1)/2\). Majority votes, list concatenation, clipping, floored
averages, and coordinate extrema take polynomially many operations
on the bounded-size values above.

The work attached to one main-chain node includes at most
\(J(g+1)\) orientation invocations and polynomially many local
vector-loop steps. Terminal \(\Basic\) calls in these orientation
invocations contribute at most
\(J(g+1)(1+128ns)\) operator evaluations; they are not additional
solver invocations in \(N_d\). Keeping every center iterate and
every exceptional-list entry still uses polynomial local data and
work. Consequently the local work attached to one node, leaving its
solver calls unexpanded, is bounded by a polynomial \(P(L)\) of
absolute degree.

\paragraph{Generating bits and implementing the recursion.}
The only random primitive selects an iterate index. Because \(s\)
is a power of two, exactly \(\log_2s\) fair bits select a uniform
index from \(\{0,\ldots,s-1\}\), without a rejection loop. The
\(J\) trials in either amplification are run sequentially, taking
fresh bits from one tape. By Proposition~\ref{prop:termination},
each preceding subroutine terminates on every tape. Under fair
randomness the unread tail after this finite stopping time supplies
independent fresh bits, as required by the correctness proof. On a
fixed arbitrary tape the same execution respects every deterministic
call bound used above. Each orientation invocation reads at most
\(\log_2s\) bits of its own, already covered by its polynomial
local cost.

All routines are finite binary computations: \(\Ask\) performs its
\(J\) solver calls, and \(\OrientPlus\) performs its \(J\)
orientation calls. The final winning set requires \(n\) comparisons
with \(B_*/2\). Enlarge \(P\) to include this output step and the
preprocessing. For the total bit cost \(T_d\), contracting the main
chain as before gives
\begin{equation}\label{cost:bit-recurrence}
 T_0\le P(L),\qquad
 T_d\le KP(L)+KA_0T_{d-1}\quad(d\ge1).
\end{equation}
The estimates used in Equation~\eqref{cost:call-count} now yield
\(T_{D_*}=2^{O((\log_2(L+2))^2)}\).

A uniform machine can perform every instruction on binary strings.
Polynomial simulation overhead for data access or recursion in a
chosen standard model changes only the absolute constant in the
exponent. Arithmetic, storage manipulation, preprocessing, random
bits, and all recursive simulations are thereby included, on every
tape, in the asserted bound.
\end{proof}
